\documentclass[11pt,a4paper]{article}

\usepackage{geometry}
\geometry{a4paper, top=1.8cm, bottom=1.8cm, left=2.0cm, right=2.0cm}

\usepackage{fontspec}
\setmainfont{DejaVu Serif}
\setsansfont{DejaVu Sans}
\setmonofont{DejaVu Sans Mono}[Scale=0.87]

\usepackage{amsmath,amssymb}
\usepackage{graphicx}
\graphicspath{{figures/}}

\usepackage{xcolor}
\usepackage{booktabs}
\usepackage{tabularx}
\usepackage{caption}
\usepackage{listings}
\usepackage{fancyhdr}
\usepackage{microtype}
\usepackage{tcolorbox}
\tcbuselibrary{skins,breakable}

\usepackage{hyperref}
\hypersetup{
    colorlinks=true,
    linkcolor=blue!75!black,
    urlcolor=teal!70!black,
    citecolor=blue!75!black,
    pdfauthor={Παναγόπουλος Αθανάσιος},
    pdftitle={Λειτουργικά Συστήματα - Δραστηριότητα 2: Χρονοπρογραμματισμός Διεργασιών στο xv6}
}

\emergencystretch=3em

% --- Listing Styles ---
\definecolor{codebg}{rgb}{0.97,0.97,0.97}
\definecolor{codeborder}{rgb}{0.85,0.85,0.85}
\definecolor{codegreen}{rgb}{0.1,0.5,0.2}
\definecolor{codegray}{rgb}{0.4,0.4,0.4}
\definecolor{codepurple}{rgb}{0.5,0.1,0.7}
\definecolor{codeblue}{rgb}{0.1,0.2,0.8}

\lstdefinestyle{cstyle}{
    backgroundcolor=\color{codebg},
    commentstyle=\color{codegreen}\itshape,
    keywordstyle=\color{codeblue}\bfseries,
    numberstyle=\tiny\color{codegray},
    stringstyle=\color{codepurple},
    basicstyle=\ttfamily\footnotesize,
    breakatwhitespace=false,
    breaklines=true,
    captionpos=b,
    keepspaces=true,
    numbers=left,
    numbersep=7pt,
    showspaces=false,
    showstringspaces=false,
    showtabs=false,
    tabsize=4,
    frame=single,
    aboveskip=3pt,
    belowskip=2pt
}

\lstdefinestyle{termstyle}{
    backgroundcolor=\color{black!92},
    basicstyle=\ttfamily\footnotesize\color{white},
    breaklines=true,
    frame=single,
    rulecolor=\color{black},
    numbers=none,
    aboveskip=3pt,
    belowskip=2pt
}

\lstset{style=cstyle}

\usepackage{enumitem}
\setlist{topsep=1.5pt, itemsep=1pt, parsep=0pt}

% --- Header / Footer ---
\pagestyle{fancy}
\fancyhf{}
\fancyhead[L]{\small\textsc{Λειτουργικά Συστήματα (ECE ΓΚ702)}}
\fancyhead[R]{\small\textsc{Δραστηριότητα 2}}
\fancyfoot[C]{\thepage}
\renewcommand{\headrulewidth}{0.4pt}

\begin{document}

% --- Header Block ---
\begin{center}
    {\large \textbf{ΠΑΝΕΠΙΣΤΗΜΙΟ ΠΑΤΡΩΝ}}\\[1mm]
    {\normalsize Τμήμα Ηλεκτρολόγων Μηχανικών \& Τεχνολογίας Υπολογιστών}\\[3mm]
    {\large \textbf{Λειτουργικά Συστήματα (ECE ΓΚ702)}}\\[1mm]
    \rule{\textwidth}{0.8pt}\\[3mm]
    {\LARGE \textbf{Δραστηριότητα 2: Χρονοπρογραμματισμός Διεργασιών στο xv6}}\\[2mm]
    \rule{\textwidth}{0.8pt}\\[3mm]
    \begin{tabular}{ll}
        \textbf{Φοιτητής:} & Παναγόπουλος Αθανάσιος \\
        \textbf{Αριθμός Μητρώου (ΑΜ):} & 1107645 \\
        \textbf{Ημερομηνία:} & 7 Νοεμβρίου 2026 \\
        \textbf{Repository:} & \url{https://github.com/RealThanosP/thanOS-v6} \\
    \end{tabular}
\end{center}
\vspace{2mm}


% ==============================================================================
% ΑΣΚΗΣΗ 1
% ==============================================================================
\section{Άσκηση 1: Αναπαράσταση διεργασίας στο xv6}

Για κάθε διεργασία το kernel κρατά μία εγγραφή τύπου \texttt{struct proc} στον πίνακα \texttt{ptable.proc}. Βρίσκουμε τον ορισμό στο \texttt{include/proc.h}.

\subsection*{Η δομή στο \texttt{include/proc.h}}
% include/proc.h:93-112
\begin{lstlisting}[language=C,firstnumber=93,basicstyle=\ttfamily\scriptsize,lineskip=-0.7pt]
struct proc
{
  uintp sz;                   // Size of process memory (bytes)
  pde_t *pgdir;               // Page table
  char *kstack;               // Bottom of kernel stack for this process
  enum procstate state;       // Process state
  volatile int pid;           // Process ID
  struct proc *parent;        // Parent process
  struct trapframe *tf;       // Trap frame for current syscall
  struct context *context;    // swtch() here to run process
  void *chan;                 // If non-zero, sleeping on chan
  int killed;                 // If non-zero, have been killed
  struct file *ofile[NOFILE]; // Open files
  struct inode *cwd;          // Current directory
  char name[16];              // Process name (debugging)

  int inuse;   // If it's being run by a CPU or not
  int ticks;   // How many ticks has accumulated
  int tickets; // Number of tickets this process has (for lottery scheduling)
};
\end{lstlisting}

\subsection*{Μνήμη και αποθηκευμένοι καταχωρητές}
\begingroup\small
\begin{tabularx}{\textwidth}{@{}l X@{}}
\toprule
\textbf{Πεδίο} & \textbf{Τι κρατάει το kernel} \\
\midrule
\texttt{sz} & Το μέγεθος της μνήμης χρήστη σε bytes. Ο τύπος \texttt{uintp} είναι ακέραιος με το πλάτος ενός δείκτη. \\
\texttt{pgdir} & Δείκτη στον πίνακα σελίδων. Μέσω αυτού η CPU μεταφράζει τις εικονικές διευθύνσεις της διεργασίας σε φυσικές. \\
\texttt{kstack} & Τη βάση της στοίβας που χρησιμοποιεί η διεργασία όταν εκτελεί κώδικα του πυρήνα. \\
\texttt{tf} & Δείκτη στο trap frame. Εκεί αποθηκεύονται οι καταχωρητές και η διεύθυνση επιστροφής όταν, για παράδειγμα, ένα system call μάς φέρει από το user space στον πυρήνα. \\
\texttt{context} & Δείκτη στο αποθηκευμένο kernel context. Από εδώ η \texttt{swtch()} βρίσκει τη στοίβα για να συνεχιστεί η εκτέλεση μέσα στον πυρήνα. \\
\bottomrule
\end{tabularx}
\endgroup

\clearpage
\subsection*{Κατάσταση, πόροι και στατιστικά της διεργασίας}
\begingroup\small
\begin{tabularx}{\textwidth}{@{}l X@{}}
\toprule
\textbf{Πεδίο} & \textbf{Τι κρατάει το kernel} \\
\midrule
\texttt{state} & Την κατάσταση της θέσης ή της διεργασίας. Μια θέση ξεκινά ως \texttt{UNUSED} και δεσμεύεται ως \texttt{EMBRYO}. Η διεργασία είναι \texttt{RUNNABLE} όσο περιμένει CPU, \texttt{RUNNING} όσο εκτελείται και \texttt{SLEEPING} όσο περιμένει κάποιο γεγονός. Με την \texttt{exit()} μένει \texttt{ZOMBIE} ώσπου η \texttt{wait()} να ελευθερώσει τη θέση. \\
\texttt{pid} & Τον αριθμό που αναγνωρίζει τη συγκεκριμένη διεργασία. \\
\texttt{parent} & Δείκτη στην εγγραφή του γονέα. Η \texttt{wait()} τον ελέγχει όταν ψάχνει ποιες διεργασίες είναι παιδιά της καλούσας. \\
\texttt{chan} & Το κανάλι αναμονής της \texttt{sleep()}. Η \texttt{wakeup(chan)} κάνει \texttt{RUNNABLE} όσες διεργασίες κοιμούνται στο ίδιο κανάλι. \\
\texttt{killed} & Σημαία που θέτει η \texttt{kill()}. το kernel την ελέγχει πριν επιστρέψει στο user space και, αν έχει τεθεί, καλεί την \texttt{exit()}. \\
\midrule
\texttt{ofile[NOFILE]} & Δείκτες στα ανοιχτά αρχεία. Ο file descriptor είναι ο αριθμός θέσης στον πίνακα, π.χ. το \texttt{ofile[1]} για την τυπική έξοδο. \\
\texttt{cwd} & Δείκτη στο inode του τρέχοντος καταλόγου, από όπου ξεκινά η αναζήτηση ενός σχετικού path. \\
\texttt{name[16]} & Το όνομα που εμφανίζει η \texttt{procdump()} για το debugging. \\
\midrule
\texttt{inuse} & Γίνεται 1 στην πρώτη εκτέλεση και 0 στην \texttt{exit()}. Παραμένει 1 και στον ύπνο. Για να δούμε αν τρέχει τώρα, ελέγχουμε το \texttt{state}. \\
\texttt{ticks} & Timer ticks σε user mode (μόνο στη CPU 0). \\
\texttt{tickets} & Τα εισιτήρια της διεργασίας. Ο αρχικός scheduler αγνοεί το πεδίο. Στην Άσκηση 3 τα χρησιμοποιούμε για την κλήρωση. \\
\bottomrule
\end{tabularx}
\endgroup

\clearpage
% ==============================================================================
% ΑΣΚΗΣΗ 2
% ==============================================================================
\section{Άσκηση 2: Οι συναρτήσεις \texttt{scheduler()} και \texttt{sched()}}

\subsection*{1. Πώς παίρνει τη CPU μια διεργασία}
Κάθε CPU τρέχει τη δική της \texttt{scheduler()} στο \texttt{kernel/proc.c}. Στην έκδοση πριν από την Άσκηση 3, η συνάρτηση σαρώνει επανειλημμένα τον \texttt{ptable.proc} με τη σειρά. Ενεργοποιεί τα interrupts με \texttt{sti()} και παίρνει το \texttt{ptable.lock} πριν εξετάσει τις εγγραφές. Έτσι, δύο CPU δεν μπορούν να επιλέξουν ταυτόχρονα την ίδια \texttt{RUNNABLE} διεργασία.

% 337f603:kernel/proc.c:265-303
\begin{lstlisting}[language=C,firstnumber=265,basicstyle=\ttfamily\scriptsize,lineskip=-0.7pt]
void scheduler(void)
{
  struct proc *p = 0;

  for (;;)
  {
    // Enable interrupts on this processor.
    sti();

    // no runnable processes? (did we hit the end of the table last time?)
    // if so, wait for irq before trying again.
    if (p == &ptable.proc[NPROC])
      hlt();

    // Loop over process table looking for process to run.
    acquire(&ptable.lock);
    for (p = ptable.proc; p < &ptable.proc[NPROC]; p++)
    {
      if (p->state != RUNNABLE)
        continue;

      // Switch to chosen process.  It is the process's job
      // to release ptable.lock and then reacquire it
      // before jumping back to us.
      proc = p;
      switchuvm(p);
      p->state = RUNNING;
      p->inuse = 1;

      swtch(&cpu->scheduler, proc->context);

      switchkvm();
      // Process is done running for now.
      // It should have changed its p->state before coming back.
      proc = 0;
    }
    release(&ptable.lock);
  }
}
\end{lstlisting}

Ας πάρουμε την περίπτωση που ο scheduler βρίσκει το shell έτοιμο. Βάζει τον δείκτη του στο \texttt{proc}, που είναι τοπικό στη συγκεκριμένη CPU, και η \texttt{switchuvm(p)} φορτώνει τον πίνακα σελίδων του shell. Στη γραμμή 294 η \texttt{swtch()} αποθηκεύει το context του scheduler στο \texttt{cpu->scheduler} και φορτώνει εκείνο του shell. Από εδώ και πέρα η CPU εκτελεί κώδικα της διεργασίας.

Η \texttt{switchkvm()} στη γραμμή 296 εκτελείται όταν ο έλεγχος γυρίσει στον scheduler. Επαναφέρει τον πίνακα σελίδων του πυρήνα. Το \texttt{proc = 0} δηλώνει ότι η συγκεκριμένη CPU βρίσκεται πάλι στον scheduler. Η εγγραφή του shell παραμένει στον \texttt{ptable.proc} με την κατάσταση που έθεσε πριν παραδώσει τη CPU. Ο αρχικός scheduler συνεχίζει τη σάρωση από την επόμενη θέση. Στο τέλος του πίνακα αφήνει το lock και περιμένει το επόμενο interrupt με \texttt{hlt()}.

\clearpage
\subsection*{2. Πώς επιστρέφει ο έλεγχος στον scheduler}
Η τρέχουσα διεργασία καλεί τη \texttt{sched()} μέσα από τον πυρήνα. Για παράδειγμα, όταν λήξει το κβάντο της, ο timer interrupt οδηγεί στη \texttt{yield()}. Αυτή παίρνει το \texttt{ptable.lock}, θέτει \texttt{state = RUNNABLE} και καλεί τη \texttt{sched()}. Το ίδιο σημείο χρησιμοποιούν η \texttt{sleep()} αφού θέσει \texttt{SLEEPING} και η \texttt{exit()} αφού θέσει \texttt{ZOMBIE}.

% 337f603:kernel/proc.c:307-322
\begin{lstlisting}[language=C,firstnumber=307,basicstyle=\ttfamily\scriptsize,lineskip=-0.7pt]
void sched(void)
{
  int intena;

  if (!holding(&ptable.lock))
    panic("sched ptable.lock");
  if (cpu->ncli != 1)
    panic("sched locks");
  if (proc->state == RUNNING)
    panic("sched running");
  if (readeflags() & FL_IF)
    panic("sched interruptible");
  intena = cpu->intena;
  swtch(&proc->context, cpu->scheduler);
  cpu->intena = intena;
}
\end{lstlisting}

Πριν από την αλλαγή, η \texttt{sched()} ελέγχει ότι κρατάμε μόνο το \texttt{ptable.lock} (\texttt{ncli == 1}), ότι τα interrupts είναι απενεργοποιημένα και ότι η διεργασία έχει ήδη φύγει από το \texttt{RUNNING}. Στη γραμμή 320 η \texttt{swtch()} σώζει καταχωρητές στη στοίβα πυρήνα της διεργασίας, γράφει τη θέση της στοίβας στο \texttt{proc->context} και φορτώνει τη στοίβα του scheduler. Το \texttt{ret} συνεχίζει από τη διεύθυνση επιστροφής που βρίσκεται σε εκείνη τη στοίβα.

Αν το shell είχε καλέσει \texttt{yield()}, την επόμενη φορά που θα επιλεγεί συνεχίζει ακριβώς μετά τη δική του \texttt{swtch()}. Η \texttt{sched()} επαναφέρει το \texttt{intena}, δηλαδή την αποθηκευμένη ένδειξη για την επανενεργοποίηση των interrupts, και η \texttt{yield()} αφήνει το lock.

\subsection*{3. Η υλοποίηση της \texttt{swtch()}}
Ο πηγαίος κώδικας του \texttt{kernel/swtch.S} δείχνει το πραγματικό πέρασμα από τη μία στοίβα πυρήνα στην άλλη. Τα \texttt{pushl} αποθηκεύουν τους καταχωρητές στην παλιά στοίβα, το \texttt{movl} αλλάζει τον \texttt{esp}, και τα \texttt{popl} μαζί με το \texttt{ret} συνεχίζουν από το αποθηκευμένο context. Αυτή είναι η παραλλαγή 32-bit· το προεπιλεγμένο 64-bit build χρησιμοποιεί το αντίστοιχο \texttt{kernel/swtch64.S}.

\lstinputlisting[basicstyle=\ttfamily\scriptsize,lineskip=-0.7pt]{../../kernel/swtch.S}

\clearpage
% ==============================================================================
% ΑΣΚΗΣΗ 3
% ==============================================================================
\section{Άσκηση 3: Lottery Scheduling στο xv6}

Με τη νέα \texttt{scheduler()} κάθε \texttt{RUNNABLE} διεργασία συμμετέχει στην κλήρωση με τα δικά της εισιτήρια. Ας πάρουμε τρεις διεργασίες με 1, 2 και 3 tickets. Κληρώνουμε έναν αριθμό από το 0 έως το 5. Το 0 ανήκει στην πρώτη, τα 1 και 2 στη δεύτερη και τα 3, 4 και 5 στην τρίτη. Άρα η τρίτη έχει πιθανότητα $3/6$ να πάρει την επόμενη σειρά στη CPU. Αν κοιμηθεί, μένουν τρία εισιτήρια και η πιθανότητα της πρώτης ανεβαίνει από $1/6$ σε $1/3$.

\subsection*{1. Αρχικά tickets, κληρονόμηση και \texttt{settickets()}}
Στο \texttt{kernel/proc.c}, η \texttt{allocproc()} δίνει ένα εισιτήριο σε κάθε νέα θέση. Έτσι και η \texttt{init} μπορεί να κερδίσει την πρώτη κλήρωση.

% kernel/proc.c:42-47
\begin{lstlisting}[language=C,firstnumber=42,basicstyle=\ttfamily\scriptsize,lineskip=-0.7pt]
found:
  p->state = EMBRYO;
  p->pid = nextpid++;
  p->ticks = 0;
  p->tickets = 1;
  release(&ptable.lock);
\end{lstlisting}
Στη \texttt{fork()} αντιγράφουμε τα tickets του γονέα πριν γίνει το παιδί \texttt{RUNNABLE}.

% kernel/proc.c:154-160
\begin{lstlisting}[language=C,firstnumber=154,basicstyle=\ttfamily\scriptsize,lineskip=-0.7pt]
  pid = np->pid;
  safestrcpy(np->name, proc->name, sizeof(proc->name));
  acquire(&ptable.lock);
  np->tickets = proc->tickets;
  np->state = RUNNABLE;
  release(&ptable.lock);
  return pid;
\end{lstlisting}
Η δήλωση στο \texttt{include/user.h}, ο αριθμός της κλήσης, το assembly wrapper και η εγγραφή στον πίνακα \texttt{syscalls[]} υπήρχαν ήδη. Αφαιρέσαμε από το \texttt{kernel/syscall.c} την προσωρινή \texttt{sys\_settickets()} που επέστρεφε πάντα $-1$. Στο \texttt{kernel/sysproc.c} η υλοποίηση ελέγχει το όρισμα και αλλάζει μόνο τα tickets της διεργασίας που έκανε την κλήση. Το 0 δηλώνει επιτυχία.

% kernel/sysproc.c:38-47
\begin{lstlisting}[language=C,firstnumber=38,basicstyle=\ttfamily\scriptsize,lineskip=-0.7pt]
int sys_settickets(void)
{
  int n;
  if (argint(0, &n) < 0 || n < 1)
    return -1;
  acquire(&ptable.lock);
  proc->tickets = n;
  release(&ptable.lock);
  return 0;
}
\end{lstlisting}

\clearpage
\subsection*{2. Η κλήρωση στον \texttt{kernel/proc.c}}
Αθροίζουμε τα εισιτήρια των \texttt{RUNNABLE} διεργασιών. Η \texttt{draw\_ticket()} κρατά όσα bits χρειάζονται και ξανατραβά αριθμό αν αυτός πέσει έξω από το διάστημα. Για έξι εισιτήρια, για παράδειγμα, η μάσκα κρατά τρία bits και οι τιμές 6 και 7 απορρίπτονται. 
Στο δεύτερο πέρασμα αφαιρούμε τα εισιτήρια κάθε διεργασίας ώσπου ο αριθμός να πέσει στο διάστημα του νικητή. Το \texttt{ptable.lock} καλύπτει και τα δύο περάσματα. Όταν η διεργασία επιστρέψει στον scheduler, το \texttt{break} οδηγεί σε καινούρια κλήρωση.

% kernel/proc.c:281-326
\begin{lstlisting}[language=C,firstnumber=281,basicstyle=\ttfamily\scriptsize,lineskip=-0.7pt]
void scheduler(void)
{
  struct proc *p;

  for (;;) {
    sti();
    acquire(&ptable.lock);
    unsigned long long total = 0;
    for (p = ptable.proc; p < &ptable.proc[NPROC]; p++)
      if (p->state == RUNNABLE)
        total += p->tickets;

    if (total == 0) {
      release(&ptable.lock);
      hlt();
      continue;
    }

    unsigned long long ticket = draw_ticket(total);
    for (p = ptable.proc; p < &ptable.proc[NPROC]; p++) {
      if (p->state != RUNNABLE)
        continue;
      if (ticket >= p->tickets) {
        ticket -= p->tickets;
        continue;
      }

      // Switch to chosen process.  It is the process's job
      // to release ptable.lock and then reacquire it
      // before jumping back to us.
      proc = p;
      switchuvm(p);
      p->state = RUNNING;
      p->inuse = 1;

      swtch(&cpu->scheduler, proc->context);

      switchkvm();
      // Process is done running for now.
      // It should have changed its p->state before coming back.
      proc = 0;
      break; // Draw again for the next dispatch.
    }
    release(&ptable.lock);
  }
}
\end{lstlisting}

\subsection*{3. Γεννήτρια}
Στο δικό μας xv6 η \texttt{rand()} βρίσκεται στο τέλος του \texttt{kernel/proc.c}, από την Δραστηριότητα 1. Την επεκτείναμε ώστε να επιστρέφει 64-bit αριθμό: το άθροισμα των tickets μπορεί να υπερβεί τα $2^{32}$, παρότι κάθε διεργασία έχει 32-bit πλήθος. Ο scheduler την καλεί μέσω της \texttt{draw\_ticket()} όσο κρατά το \texttt{ptable.lock}.

\clearpage
\subsection*{4. Κώδικας του \texttt{user/schedtest.c}}
Ο γονέας δημιουργεί τρία παιδιά που εκτελούν ατέρμονα βρόχο. Διαβάζει τους μετρητές τους πριν και μετά από 600 ticks, τυπώνει τη διαφορά και τα τερματίζει.
Η δοκιμή με $2^{31}-1$ tickets ανά παιδί βασίζεται στη νέα 64-bit \texttt{rand()}, ώστε η κλήρωση να καλύπτει όλο το άθροισμα των tickets.

\lstinputlisting[language=C,basicstyle=\ttfamily\scriptsize,lineskip=-1.5pt]{../../user/schedtest.c}

\clearpage
\subsection*{5. Δοκιμή με τρεις διεργασίες σε μία CPU}
Ξεκινάμε το QEMU με \texttt{make CPUS=1 qemu-nox}. Στο shell του xv6 δίνουμε, για παράδειγμα, \texttt{schedtest 1 2 3}. Πριν από κάθε \texttt{fork()}, ο γονέας καλεί τη \texttt{settickets(n)}. Με τη \texttt{getpinfo()} διαβάζει τους αρχικούς μετρητές, κοιμάται για 600 ticks και τους ξαναδιαβάζει. Η διαφορά δίνει τον χρόνο CPU που πήρε κάθε παιδί στο ίδιο διάστημα.

\begin{center}
\begin{tabular}{@{}lrrr@{}}
\toprule
Tickets ανά παιδί & Ticks 1ου & Ticks 2ου & Ticks 3ου \\
\midrule
$1,\ 1,\ 1$ & 197 & 209 & 194 \\
$1,\ 2,\ 3$ & 103 & 196 & 301 \\
$3,\ 2,\ 1$ & 285 & 194 & 121 \\
$2^{31}-1$ στο καθένα & 195 & 205 & 203 \\
\bottomrule
\end{tabular}
\end{center}

Με ίσα tickets περιμένουμε περίπου 200 ticks ανά παιδί. 
Με την αναλογία $1:2:3$ περιμένουμε 100, 200 και 300. 
Το τρίτο παιδί πήρε 301 από τα 600 ticks ($50{,}2\%$). 
Όταν αντιστρέψαμε τα tickets, το πρώτο παιδί πήρε τα περισσότερα (285). 
Οι μικρές αποκλίσεις είναι αναμενόμενες, αφού σε κάθε επιλογή γίνεται νέα κλήρωση. 
Η τελευταία δοκιμή ελέγχει και το άθροισμα. 
Τα τρία παιδιά έχουν μαζί $6\,442\,450\,941$ tickets, οπότε χρειαζόμαστε περισσότερα από 32 bits.
Τα αποτελέσματα διαφέρουν, από run σε run του \texttt{schedtest} λόγω του randomization του lottery scheduling, οπότε τα αποτελέσματα είναι ενδεικτικά.
\end{document}
